\documentclass[brazilian,a4paper]{beamer}

\usetheme{CambridgeUS}
\usecolortheme{beaver}

% \usepackage[utf8]{inputenc}
\usepackage{babel}
\usepackage{ragged2e}
\apptocmd{\frame}{}{\justifying}{}

\usepackage{xcolor}
\definecolor{emerald}{rgb}{0.31, 0.78, 0.47}

\usepackage{caption}
\usepackage{hyperref}
\hypersetup{
    unicode=true,
    pdfstartview={FitV},
    pdfauthor={Nightwind},
    pdfkeywords={latex, dicas},
    colorlinks=true,
    linkcolor=black,
    citecolor=green,
    filecolor=cyan,
    urlcolor=magenta
}
\usepackage[nameinlink,capitalise]{cleveref}
\usepackage[style=abnt]{biblatex}
\addbibresource{references.bib}
\usepackage{csquotes}

\usepackage{tabularx}
\usepackage{array}
\newcolumntype{M}[1]{>{\raggedleft\arraybackslash}m{#1}}
\newcolumntype{O}[1]{>{\raggedright\arraybackslash}m{#1}}
\newcolumntype{C}[1]{>{\centering\arraybackslash}m{#1}}
\usepackage{multirow}

\usepackage{graphicx}
\usepackage{svg}

\usepackage{multicol}
\usepackage{fontspec}
\setsansfont{Syne}
\setmonofont{Space Mono}
\usefonttheme[onlymath]{serif}

\usepackage{mathtools}
\usepackage{mathrsfs}
\usepackage{xfrac}
\DeclareMathOperator{\sen}{sen}
\newcommand{\intd}[4]{\ensuremath{\int_{#1}^{#2}\left[#3\right]\,\mathsf{d}{#4}}}
\newcommand{\edp}[3][]{\ensuremath{\frac{\partial^{#1}\left[#2\right]}{\partial{#3}^{#1}}}}

\usepackage{listings}
\definecolor{asparagus}{rgb}{0.53, 0.66, 0.42}
\definecolor{lavendergray}{rgb}{0.77, 0.76, 0.82}
\definecolor{pastelmagenta}{rgb}{0.96, 0.6, 0.76}
\definecolor{mistyrose}{rgb}{1.0, 0.89, 0.88}
\definecolor{packagecolor}{rgb}{1.0, 0.0, 0.16}
\definecolor{darkslateblue}{rgb}{0.28, 0.24, 0.55}
\definecolor{forestgreen}{rgb}{0.13, 0.55, 0.13}

\lstdefinestyle{myStyleLatex}{
    language=[LaTeX]TeX,
    backgroundcolor = \color{mistyrose},
    basicstyle = \ttfamily,
    breakatwhitespace = true,
    columns = fullflexible,
    breaklines = true,
    captionpos = a,
    commentstyle = {\footnotesize\color{asparagus}},
    escapeinside = {  {(*}  {*)}  },
    extendedchars = true,
    firstnumber = 1,
    frame = none,
    keepspaces = true,
    keywordstyle = {\bfseries\color{darkslateblue}},
    keywordstyle = [2]{\color{forestgreen}},
    keywordstyle = [3]{\itshape\color{packagecolor}},
    morekeywords = {maketitle,chapter,part, section, subsection, subsubsection, paragraph, subparagraph,familydefault,rmdefault,sfdefault,ttdefault,textsubscript,ttshape,colorbox,textcolor,definecolor,setlength,includegraphics,listoffigures,listoftables,endfirsthead,endhead,endfoot,endlastfoot,arraybackslash,newcolumntype,rowcolor,rowcolors,cellcolor,multirow,setmainfont,documentclass,mathcal,mathbb,mathfrak,mathscr,DeclareMathOperator,cref,autoref,lstdefinestyle,lstset,dfrac,sfrac,lstinline,lstinputlisting},
    morekeywords = [2]{default,arguments},
    morekeywords = [3]{document,command,definition, ambiente,book,report,article,exam, beamer,flushright,flushleft,center,figure,table,tabular,tabularx,longtable,article,book,exam,wrapfigure,equation,split,align,gather,multiline,lstlisting},
    numbers = left,
    numbersep = -1pt,
    numberstyle = \color{lavendergray},
    rulecolor = \color{magenta},
    showspaces = false,
    showstringspaces = false,
    showtabs = false,
    stepnumber = 1,
    stringstyle = \color{pastelmagenta},
    tabsize = 2
}

\usepackage[os=win]{menukeys}

\author{Nightwind}
\institute[CTISM]{Colégio Técnico Industrial de Santa Maria}
\logo{\includegraphics[width=1cm]{../images/photo1.jpg}}
\date{\today}


\title{Citação e Referências pela ABNT}

\begin{document}
\frame{\titlepage}

\begin{frame}
    \frametitle{Sumário}
    \tableofcontents
\end{frame}

\section{Introdução}
\begin{frame}
    \frametitle{Introdução}

    \begin{itemize}
        \item As referências são as partes mais cruciais da escrita científica. Elas garantem a credibilidade e o conhecimento sobre o assunto. 
        \item Normalmente, por mais que a cobrança de referenciação seja recorrente, a formatação varia de acordo com a instituição ou revista que receberá o material.
        \item Essa formatação é complexa de se fazer manualmente: cada origem de referência possui um padrão e cada instituição cobra um padrão diferente.
        \item Para o \LaTeX, esse é um ponto que apresenta maior facilidade do que nos editores de texto normais.
        \item Insere-se a referência no formato \texttt{.bib}, após a primeira compilação, pode citar através de um comando simples com o rótulo da referência e, onde se desejar exibir a listagem de referências é só adicionar outro comando.
    \end{itemize}

\end{frame}

\section{Pacote}
\begin{frame}[fragile]
    \frametitle{Pacote}

    \begin{itemize}
        \item O pacote sugerido é o \texttt{biblatex}. Ele vai garantir que o estilo abnt seja utilizado. 
        \item \textbf{Atentar para a versão de ABNT requerida e a versão disponível.}
        \item É recomendável utilizar o pacote \texttt{babel} para o suporte às diferentes linguagens. E o pacote \texttt{csquotes} de acordo com a linguagens (adaptada pelo \texttt{babel}) ajusta melhor as demarcações de citações de aspas, por exemplo.
    \end{itemize}

\end{frame}

\subsection{Opções}
\begin{frame}[fragile]
    \frametitle{Opções}

    As opções de para configuração do pacote são:
    \begin{itemize}
        \item \texttt{style}: \texttt{numeric}, \texttt{alphabetic}, \ldots, \texttt{\bfseries abnt}.
        \item \texttt{backend}: considerando que termos acentos e cedilha, sugiro usar o \texttt{biber}.
        \item Acessar o manual para mais opções: \citeurl{biblatexPackage}.
    \end{itemize}

\end{frame}

\section{Arquivo de Bibliografia}
\begin{frame}[fragile,allowframebreaks]
    \frametitle{Arquivo de Bibliografia}

    \begin{itemize}
        \item O arquivo de bibliografia possui uma extensão diferente. Ele é \texttt{.bib}
        \item A formatação dele também é diferente. Exemplo:
    \end{itemize}
\begin{lstlisting}[style=myStyleLatex]
@article{<nome>,
    title={},
    author={},
    journal={},
    volume={},
    number={},
    pages={},
    year={},
    publisher={}
}

@book{<nome>,
    author = {},
    year = {},
    title = {},
    publisher = {},
    address = {},
    edition = {}
}
\end{lstlisting}
\end{frame}

\begin{frame}[fragile]
    \frametitle{Arquivo de Bibliografia}

    \begin{itemize}
        \item Para conectar o arquivo de referências ao arquivo \texttt{.tex}, colocar no preâmbulo o seguinte comando:
    \end{itemize}
\begin{lstlisting}[style=myStyleLatex]
\addbibresource{<file.bib>}
\end{lstlisting}

    Pode-se adicionar mais de um documento, chamando o código inteiro novamente. Atentar para não ter duas entradas com o mesmo nome.

\end{frame}


\section{Citação}
\begin{frame}[fragile]
    \frametitle{Citar um Referência}

    Para citar uma referência, recomenda-se que o arquivo \texttt{bib} esteja salvo e que o arquivo já tenha passado por, no mínimo, uma compilação. Desse modo, os comandos abaixo reconhece a lista dos indicadores. 

\begin{lstlisting}[style=myStyleLatex]
\cite[<info>]{<nome>}
\cites[<info>]{<nome>}[<info>]{<nome>}
\textcite[<info>]{<nome>}
\end{lstlisting}

\end{frame}

\begin{frame}[fragile]
    \frametitle{Não citar Referência}

    \begin{itemize}
        \item Usar o comando \lstinline[style=myStyleLatex]!\nocite{}!, nele pode ter um asterisco para chamar todas as entradas do arquivo \texttt{bib}. Ou todas as entradas que se objetiva demonstrar no documento, mas que não se faz a citação direta ou indireta.
    \end{itemize}

\end{frame}

\section{Referências}

\begin{frame}[fragile]
    \frametitle{Referências}

    \begin{itemize}
        \item Para apresentar as referências. Usar o comando \lstinline[style=myStyleLatex]!\printbibliography!
    \end{itemize}

\end{frame}


\begin{frame}[allowframebreaks]
    \frametitle{Referências}

    \nocite{*}
    \printbibliography[keyword=inserirReferencias]

\end{frame}

\end{document}